% !PRINTED 475
\weberpage{475}
%<<CONT
und benutzen die aus § 34, (19) folgende Formel:
\[
f(\omega)\,f_2\left(\frac{-3+\omega}{2}\right) = e^{-\frac{\pi i}{8}}\sqrt{2}.
\]
Demnach ist, wenn wir
\[
f(\sqrt{-m}) = x
\]
setzen, $x$ nach (1) die reelle positive Wurzel der Gleichung
\weq{5}{x^{24} + \gamma_2\left(\frac{-3+\sqrt{-m}}{2}\right) x^{16} - 2^8 = 0.}
Für $m = 7$ erhalten wir nach § 125, (18)
\weq{6}{f(\sqrt{-7}) = \sqrt{2},}
während in den anderen Fällen, nach Einsetzen der Werte für $\gamma_2$,
sich die Gleichung (5) erst in zwei Faktoren 12ten, diese wieder
in zwei Faktoren 6ten und schließlich diese in zwei 3ten Grades
spaltet. Die so erhaltenen Gleichungen sind kubische Gleichungen
für $x^8$, $x^4$, $x^2$, $x$, von denen wir jedesmal nur die beibehalten,
die eine reelle positive Wurzel hat, der schließlich $f(\sqrt{-m})$
selbst genügt.

Um die erste Zerlegung zu finden, setzen wir die Gleichung
(5) in die Form:
\[
(x^{12} - a x^4)^2 - (b x^8 + c)^2 = 0,
\]
und haben die \gesperrt{ganzen Zahlen} $a$, $b$, $c$ aus den Gleichungen:
\[
2a + b^2 = -\gamma_2, \quad 2bc = a^2, \quad c^2 = 2^8
\]
zu bestimmen, also $c = \pm 16$; die Gleichung 12ten Grades mit
positiver Wurzel lautet also:
\[
x^{12} - b x^8 - a x^4 - 16 = 0.
\]
Die Zahlen $a$, $b$ bestimmen sich aus den obigen Gleichungen
leicht, und so findet man schließlich die gesuchten kubischen
Gleichungen. Man erhält z.\,B. für $m = 11$ successive
\[
\begin{aligned}
x^{12} - 8x^8 + 16x^4 - 16 &= 0,\\
x^6 - 4x^2 - 4 &= 0,\\
x^3 - 2x^2 + 2x - 2 &= 0.
\end{aligned}
\]
In den fünf Fällen des § 125 findet man folgende Gleichungen:
\weq{7}{x = f(\sqrt{-11}), \quad x^3 - 2x^2 + 2x - 2 = 0,}
\weq{8}{x = f(\sqrt{-19}), \quad x^3 - 2x - 2 = 0,}
\weq{9}{x = f(\sqrt{-43}), \quad x^3 - 2x^2 - 2 = 0,}
\weq{10}{x = f(\sqrt{-67}), \quad x^3 - 2x^2 - 2x - 2 = 0,}
\weq{11}{x = f(\sqrt{-163}), \quad x^3 - 6x^2 + 4x - 2 = 0.}
